Calcula el \textbf{volumen}, el \textbf{área lateral} y el \textbf{área total} de las siguientes figuras:

		\begin{parts}
			\part Pirámide cuyos lados $l$  de la base miden 16 cm y la altura $h$  mide 27 cm.

			\begin{minipage}{.25\linewidth}
				\includegraphics[width=0.6\linewidth]{mex_0029.png}
			\end{minipage}\hfill
			\begin{minipage}{.75\linewidth}
				Volumen:\\[-2em]
				\begin{solutionbox}{1.5cm}
					\[V=\dfrac{1}{3}A_b h= \dfrac{1}{3}l^2h= \dfrac{1}{3}16^2(27)=2304\]
				\end{solutionbox}

				A. Lateral:\\[-2em]
				\begin{solutionbox}{1.5cm}
					\[A_L=n\dfrac{lh}{2}=4\cdot \dfrac{16\times 27}{2}=864\]
				\end{solutionbox}

				A. Total:\\[-2em]
				\begin{solutionbox}{1.5cm}
					\[A_T=A_L+l^2=864+16^2=864+256=1120\]
				\end{solutionbox}
			\end{minipage}

			\part Pirámide de 19 cm de altura cuya base es un pentágono cuyos lados  $l$  miden 8 cm y su apotema mide 5 cm.

			\begin{minipage}{.25\linewidth}
				\includegraphics[width=0.8\linewidth]{mex_0027.png}
			\end{minipage}\hfill
			\begin{minipage}{.75\linewidth}
				Volumen:\\[-2em]
				\begin{solutionbox}{1.5cm}
					\[V=\dfrac{1}{3}A_b\cdot h= \dfrac{1}{3}\left(\dfrac{nla}{2}\right)h= \dfrac{5(8)5}{2}(19)=950\]
				\end{solutionbox}
				
				A. Lateral:\\[-2em]
				\begin{solutionbox}{1.5cm}
					\[A_L=n\dfrac{lh}{2}=5\cdot 8\cdot 19=760\]
				\end{solutionbox}
				
				A. Total:\\[-2em]
				\begin{solutionbox}{1.5cm}
					\[A_T=A_L+\dfrac{nla}{2}=760+100=860\]
				\end{solutionbox}
			\end{minipage}
		\end{parts}